\documentclass[10pt,letterpaper]{article}
\usepackage{cornell-notes}

\title{Clause 17.03.05.03: Numeric Limits Specializations}
\author{}
\date{}
\setDocTitle{Clause 17.03.05.03: Numeric Limits Specializations}
\setDocSubtitle{C++ 2024}
\setDocOwner{C++ 2024 Cornell Notes}
\setDocDate{\today}
\setCornellCollection{C++ 2024 Cornell Notes}
\setCornellUnitType{Clause}
\setCornellUnitNumber{17.03.05.03}
\setCornellUnitTitle{Numeric Limits Specializations}
\hypersetup{pdftitle={Clause 17.03.05.03: Numeric Limits Specializations},pdfsubject={Cornell notes},pdfauthor={}}

\begin{document}
\maketitle
\begin{CornellOverviewBox}{Overview}
\textbf{Classification:} Standard library - Normative\\[2pt]
\textbf{Learning objective:} Understand the rules, terminology, and semantic consequences associated with numeric\_limits specializations.
\end{CornellOverviewBox}\begin{CornellSummaryBox}{Summary}
{\sffamily\bfseries\color{Accent} Summary}\par\vspace{3pt}Section 17.3.5.3 focuses on numeric\_limits specializations. Use the listed grammar productions together with the semantic constraints and disambiguation rules referenced by the section. When applying it, preserve the distinction between mandatory requirements, recommendations, permissions, and behavior deliberately left undefined, unspecified, or implementation-defined.
\end{CornellSummaryBox}

\end{document}
